\chapter{De binomiale verdeling}\label{ch:g11:binom}

Herhaal hetzelfde ja-neeexperiment een aantal keer, onafhankelijk, en tel de
successen: de verdeling die je zo krijgt --- de \emph{binomiale} --- is de
belangrijkste discrete \omterm{def:g11:prob:rv}{kansverdeling} van allemaal. Dit hoofdstuk bouwt haar op
met boomdiagrammen en het tellen van paden; de gesloten formule voor die
aantallen paden (met faculteiten) komt met het telgereedschap van
\cref{ch:g12:comb}, en de verdeling keert terug in \cref{ch:g12:randvar}.

\section{Bernoulli-experimenten}

\begin{definition}[Bernoulli-experiment]\label{def:g11:binom:bernoulli}
Een \emph{bernoulli-experiment}\index{bernoulli-experiment} is een experiment
met precies twee uitkomsten: \emph{succes}, met \omterm{def:g10:proba:distribution}{kans} $p$, en \emph{mislukking},
met \omterm{def:g10:proba:distribution}{kans} $1 - p$. Van de \omterm{def:g11:prob:rv}{toevalsvariabele} $X$ die gelijk is aan $1$ bij succes
en aan $0$ bij mislukking zegt men dat ze de \emph{bernoulli-verdeling}
$\mathcal B(p)$ volgt; dan is
\[
\E(X) = p,
\qquad
\V(X) = p(1 - p).
\]
\end{definition}

\begin{proof}[Bewijs van de twee formules]
$\E(X) = p \times 1 + (1-p) \times 0 = p$; en omdat $X^2 = X$ (zowel $0$ als
$1$ is zijn eigen kwadraat), is $\E(X^2) = p$, zodat
\cref{prop:g11:prob:konig} geeft dat $\V(X) = p - p^2 = p(1-p)$.
\end{proof}

\begin{definition}[Herhaalde onafhankelijke experimenten]\label{def:g11:binom:repeated}
Een \omterm{def:g11:binom:bernoulli}{bernoulli-experiment} $n$ keer \emph{onafhankelijk} herhalen betekent: de
uitkomst van elk experiment heeft geen invloed op de andere, en de \omterm{def:g10:proba:distribution}{kans} van
een volledige \omterm{def:g11:seq:sequence}{rij} uitkomsten is het \emph{product} van de kansen langs het
bijhorende pad van het \omterm{met:g10:proba:tree}{boomdiagram} --- $p$ voor elk succes en $1 - p$ voor
elke mislukking.
\end{definition}

\begin{example}\label{ex:g11:binom:threetrials}
Drie onafhankelijke experimenten met succeskans $p$. De \omterm{def:g11:seq:sequence}{rij} SMS (succes,
mislukking, succes) heeft \omterm{def:g10:proba:distribution}{kans} $p(1-p)p = p^2(1-p)$ --- en dat geldt voor
\emph{elke} \omterm{def:g11:seq:sequence}{rij} met precies twee successen, waar ze ook staan: alleen het
\emph{aantal} S'en en M'en telt.
\end{example}

\section{Aantallen paden en binomiaalcoëfficiënten}

\begin{definition}[Binomiaalcoëfficiënt]\label{def:g11:binom:coefficient}
In het \omterm{met:g10:proba:tree}{boomdiagram} van $n$ onafhankelijke experimenten is de
\emph{binomiaalcoëfficiënt}\index{binomiaalcoëfficiënt} $\binom{n}{k}$ (lees:
``$n$ boven $k$'') het aantal paden met precies $k$ successen.
\end{definition}

\begin{example}
$\binom{3}{2} = 3$: de paden SSM, SMS, MSS. Net zo is $\binom{3}{0} = 1$ (het
pad MMM), $\binom{3}{1} = 3$ en $\binom{3}{3} = 1$. Per afspraak, en volgens
het \omterm{met:g10:proba:tree}{boomdiagram}, is $\binom n0 = \binom nn = 1$ voor elke $n$.
\end{example}

\begin{omfigure}
\begin{tikzpicture}[font=\small]
  \coordinate (root) at (0,0);
  \node (S) at (1.5,1.5) {S};
  \node (F) at (1.5,-1.5) {M};
  \node (SS) at (3.0,2.25) {S};
  \node (SF) at (3.0,0.75) {M};
  \node (FS) at (3.0,-0.75) {S};
  \node (FF) at (3.0,-2.25) {M};
  \node (SSS) at (4.5,2.6) {S};
  \node (SSF) at (4.5,1.85) {M};
  \node (SFS) at (4.5,1.1) {S};
  \node (SFF) at (4.5,0.35) {M};
  \node (FSS) at (4.5,-0.35) {S};
  \node (FSF) at (4.5,-1.1) {M};
  \node (FFS) at (4.5,-1.85) {S};
  \node (FFF) at (4.5,-2.6) {M};
  \draw[omThm, thick] (root) -- (S);
  \draw[omThm, thick] (S) -- (SS);   \draw[omThm, thick] (SS) -- (SSF);
  \draw[omThm, thick] (S) -- (SF);   \draw[omThm, thick] (SF) -- (SFS);
  \draw[omThm, thick] (root) -- (F);
  \draw[omThm, thick] (F) -- (FS);   \draw[omThm, thick] (FS) -- (FSS);
  \draw[gray] (SS) -- (SSS);
  \draw[gray] (SF) -- (SFF);
  \draw[gray] (FS) -- (FSF);
  \draw[gray] (F) -- (FF);
  \draw[gray] (FF) -- (FFS);
  \draw[gray] (FF) -- (FFF);
  \node[gray, anchor=west] at (5.0,2.6) {$3$ successen};
  \node[omThm, anchor=west] at (5.0,1.85) {$2$ successen};
  \node[omThm, anchor=west] at (5.0,1.1) {$2$ successen};
  \node[gray, anchor=west] at (5.0,0.35) {$1$ succes};
  \node[omThm, anchor=west] at (5.0,-0.35) {$2$ successen};
  \node[gray, anchor=west] at (5.0,-1.1) {$1$ succes};
  \node[gray, anchor=west] at (5.0,-1.85) {$1$ succes};
  \node[gray, anchor=west] at (5.0,-2.6) {$0$ successen};
\end{tikzpicture}

{\small Het \omterm{met:g10:proba:tree}{boomdiagram} van $n = 3$ experimenten: $\binom{3}{2} = 3$ paden
(rood) dragen precies twee successen, elk met \omterm{def:g10:proba:distribution}{kans} $p^2(1-p)$.}
\end{omfigure}

\begin{proposition}[Regel van Pascal]\label{prop:g11:binom:pascal}
\index{driehoek van Pascal}
Voor $1 \leq k \leq n - 1$ geldt
\[
\binom{n}{k} = \binom{n-1}{k-1} + \binom{n-1}{k}.
\]
\end{proposition}

\begin{proof}
Sorteer de paden met $k$ successen in het \omterm{met:g10:proba:tree}{boomdiagram} van $n$ experimenten
volgens hun \emph{laatste} experiment. De paden die op een succes eindigen
ontstaan uit een pad van de eerste $n-1$ experimenten met $k - 1$ successen:
daarvan zijn er $\binom{n-1}{k-1}$. De paden die op een mislukking eindigen
verlengen een pad met $k$ successen onder de eerste $n-1$ experimenten:
daarvan zijn er $\binom{n-1}{k}$. Elk pad is van precies één van beide
soorten.
\end{proof}

De regel van Pascal brengt de coëfficiënten \omterm{def:g11:seq:sequence}{rij} per rij voort --- elk getal is
de som van de twee erboven:
\[
\begin{array}{ccccccccccc}
&&&&&1&&&&&\\
&&&&1&&1&&&&\\
&&&1&&2&&1&&&\\
&&1&&3&&3&&1&&\\
&1&&4&&6&&4&&1&\\
1&&5&&10&&10&&5&&1
\end{array}
\]

\begin{remark}
Een gesloten formule, $\binom nk = \frac{n!}{k!(n-k)!}$, samen met een
stelselmatige theorie van het tellen, wordt in \cref{ch:g12:comb} opgebouwd.
Op dit niveau berekent de \omterm{prop:g11:binom:pascal}{driehoek van Pascal} elke coëfficiënt die we nodig
hebben.
\end{remark}

\section{De binomiale verdeling}

\begin{theorem}[Binomiale verdeling]\label{thm:g11:binom:binomial}
Zij $X$ het aantal successen in $n$ onafhankelijke \omterm{def:g11:binom:bernoulli}{bernoulli-experimenten} met
parameter $p$. Dan volgt $X$ de \emph{binomiale verdeling}\index{binomiale
verdeling} $\mathcal B(n, p)$:
\[
\P(X = k) = \binom{n}{k}\, p^k (1-p)^{n-k},
\qquad k = 0, 1, \dots, n .
\]
\end{theorem}

\begin{proof}
De \omterm{def:g11:prob:model}{gebeurtenis} $X = k$ is de verzameling van alle paden met precies $k$
successen. Elk zo'n pad heeft \omterm{def:g10:proba:distribution}{kans} $p^k(1-p)^{n-k}$: het product langs het pad
bevat $k$ factoren $p$ en $n - k$ factoren $1-p$, in een of andere volgorde
(\cref{def:g11:binom:repeated}). Er zijn $\binom nk$ zulke paden
(\cref{def:g11:binom:coefficient}), en hun kansen tellen op.
\end{proof}

\begin{example}\label{ex:g11:binom:quiz}
Een quiz heeft $5$ onafhankelijke vragen met elk $4$ keuzemogelijkheden; een
leerling antwoordt willekeurig, zodat elke vraag een succes is met
$p = \frac14$. Het aantal $X$ juiste antwoorden volgt
$\mathcal B\left(5, \frac14\right)$, en met rij $5$ van de \omterm{prop:g11:binom:pascal}{driehoek van
Pascal}:
\[
\P(X = 2) = \binom52 \left(\frac14\right)^2\left(\frac34\right)^3
= 10 \times \frac{1}{16} \times \frac{27}{64}
= \frac{270}{1024} \approx 0.26 .
\]
De \omterm{def:g10:proba:distribution}{kans} op minstens één juist antwoord gebruikt het \omterm{def:g10:proba:operations}{complement}:
$\P(X \geq 1) = 1 - \P(X = 0) = 1 - \left(\frac34\right)^5 \approx 0.76$.
\end{example}

\begin{proposition}[Verwachtingswaarde en variantie]\label{prop:g11:binom:expectation}
Is $X \sim \mathcal B(n, p)$, dan geldt
\[
\E(X) = np,
\qquad
\V(X) = np(1-p).
\]
\end{proposition}

\begin{proof}[Verantwoording]
Schrijf $X = X_1 + X_2 + \dots + X_n$, waarbij $X_i$ gelijk is aan $1$
wanneer het $i$-de experiment lukt: elke $X_i$ is een bernoulli-variabele met
\omterm{def:g11:prob:expectation}{verwachtingswaarde} $p$ (\cref{def:g11:binom:bernoulli}). \omterm{def:g11:stat:mean}{Gemiddelden} tellen op
--- de $n$ bijdragen sommeren geeft $\E(X) = np$. Dat ook de
\emph{\omterm{def:g11:prob:variance}{varianties}} optellen voor onafhankelijke variabelen is waar maar
delicater: de formule voor de \omterm{def:g11:prob:variance}{variantie} wordt \emph{op dit niveau aangenomen}
en in \cref{ch:g12:sums} bewezen.
\end{proof}

\begin{omfigure}
\begin{tikzpicture}
\begin{axis}[
  width=10cm, height=5cm,
  ybar, bar width=9pt,
  xlabel={$k$}, ylabel={$\P(X = k)$},
  xmin=-0.8, xmax=10.8, ymin=0, ymax=0.3,
  xtick={0,1,2,3,4,5,6,7,8,9,10},
  ytick={0.1, 0.2},
  tick label style={font=\small},
  label style={font=\small},
  title={$\mathcal B(10,\ 0.5)$}]
  \addplot[fill=omDef!40, draw=omDef] coordinates {
    (0,0.00098) (1,0.00977) (2,0.04395) (3,0.11719) (4,0.20508)
    (5,0.24609) (6,0.20508) (7,0.11719) (8,0.04395) (9,0.00977)
    (10,0.00098)};
\end{axis}
\end{tikzpicture}

{\small De verdeling $\mathcal B(10, 0.5)$ (tien worpen met een eerlijke
munt): gecentreerd in $\E(X) = np = 5$, symmetrisch, met vrijwel de hele \omterm{def:g10:proba:distribution}{kans}
tussen $2$ en $8$.}
\end{omfigure}

\begin{method}[Een binomiale situatie herkennen]\label{met:g11:binom:recognize}
Ga, vóór je $X \sim \mathcal B(n, p)$ opschrijft, drie ingrediënten na: een
\emph{vast aantal} $n$ experimenten, op voorhand vastgelegd; elk experiment
heeft \emph{twee uitkomsten} met \emph{dezelfde} succeskans $p$; en de
experimenten zijn \emph{onafhankelijk} (met teruglegging, of met aparte
toestellen). Trekken zonder teruglegging uit een kleine populatie is
\emph{niet} binomiaal --- de \omterm{def:g10:proba:distribution}{kans} verandert bij elke trekking
(\cref{exo:g11:prob:6}).
\end{method}

\section{Steekproeven: is de waarneming verrassend?}

De binomiale verdeling beantwoordt een heel praktische vraag: \emph{als de
succeskans werkelijk $p$ is, welke aantallen successen zijn dan
aannemelijk?}

\begin{example}\label{ex:g11:binom:decision}
Een machine hoort hoogstens $10\%$ defecte stukken te maken. In een partij
van $10$ stukken zijn er $4$ defect. Pech of kapotte machine? Is de machine
in orde, dan volgt het aantal defecten $\mathcal B(10, 0.1)$, en is
\[
\P(X \geq 4) = 1 - \P(X \leq 3) \approx 1 - 0.987 = 0.013 :
\]
ongeveer één \omterm{def:g10:proba:distribution}{kans} op $80$. Zo'n onwaarschijnlijke \omterm{def:g11:prob:model}{gebeurtenis} waarnemen is
een sterk signaal --- men \emph{verwerpt} de veronderstelling dat de machine
nog op $10\%$ werkt, met in het achterhoofd dat de beslissing met \omterm{def:g10:proba:distribution}{kans}
ongeveer $0.013$ fout kan zijn.
\end{example}

\begin{method}[Beslissingsregel bij een binomiaal model]\label{met:g11:binom:decision}
Om een waargenomen aantal $k$ successen te beoordelen tegenover de
veronderstelling $X \sim \mathcal B(n, p)$: bereken onder die veronderstelling
de \omterm{def:g10:proba:distribution}{kans} op een resultaat dat \emph{minstens even extreem} is als $k$. Is die
\omterm{def:g10:proba:distribution}{kans} erg klein (een gangbare afspraak: onder $5\%$), verwerp dan de
veronderstelling; anders is de waarneming ermee verenigbaar. De drempel is een
keuze, geen stelling --- de statistiek kwantificeert het risico, en de
gebruiker aanvaardt het.
\end{method}

\section{Oefeningen}

\begin{exercise}[$\star$]\label{exo:g11:binom:1}
Breid de \omterm{prop:g11:binom:pascal}{driehoek van Pascal} uit tot rij $7$, en geef de waarden van
$\binom62$, $\binom{7}{3}$ en $\binom74$.
\end{exercise}

\begin{exercise}[$\star$]\label{exo:g11:binom:2}
Er wordt $4$ keer met een eerlijke dobbelsteen gegooid; $X$ telt de zessen.
Verantwoord dat $X \sim \mathcal B\left(4, \frac16\right)$ en bereken
$\P(X = 0)$, $\P(X = 1)$ en $\P(X \geq 2)$.
\end{exercise}

\begin{exercise}[$\star$]\label{exo:g11:binom:3}
Welke van de volgende situaties is binomiaal? Verantwoord.
\begin{enumerate}
  \item Het aantal keer kop bij $20$ worpen met een eerlijke munt.
  \item Het aantal azen bij $5$ kaarten uit één spel.
  \item Het aantal regendagen volgende week, als elke dag onafhankelijk met
        \omterm{def:g10:proba:distribution}{kans} $0.3$ regenachtig is.
\end{enumerate}
\end{exercise}

\begin{exercise}[$\star$]\label{exo:g11:binom:4}
$X \sim \mathcal B(50, 0.2)$. Geef $\E(X)$, $\V(X)$ en $\sigma(X)$.
\end{exercise}

\begin{exercise}[$\star\star$]\label{exo:g11:binom:5}
Een boogschutter raakt het doel bij elk schot onafhankelijk met \omterm{def:g10:proba:distribution}{kans} $0.7$.
Bereken bij $6$ schoten de \omterm{def:g10:proba:distribution}{kans} op precies $4$ treffers, en op minstens $5$
treffers.
\end{exercise}

\begin{exercise}[$\star\star$]\label{exo:g11:binom:6}
Een juist-of-fouttoets heeft $8$ vragen; een leerling gokt elk antwoord. Wat
is de \omterm{def:g10:proba:distribution}{kans} om te slagen (minstens $6$ juiste antwoorden)?
\end{exercise}

\begin{exercise}[$\star\star$]\label{exo:g11:binom:7}
Elke gekochte doos ontbijtgranen bevat figuurtje A of figuurtje B, elk
onafhankelijk met \omterm{def:g10:proba:distribution}{kans} $\frac12$. Een verzamelaar koopt $5$ dozen. Bereken de
\omterm{def:g10:proba:distribution}{kans} dat hij minstens één figuurtje van elke soort krijgt. (\omterm{def:g10:proba:operations}{Complement}: alles
A of alles B.)
\end{exercise}

\begin{exercise}[$\star\star$]\label{exo:g11:binom:8}
Een basketbalspeler scoort een vrijworp onafhankelijk met \omterm{def:g10:proba:distribution}{kans} $p$. Zij
$X \sim \mathcal B(3, p)$ het aantal rake worpen op drie pogingen. Druk
$\P(X = 3)$ en $\P(X \geq 1)$ uit als \omterm{def:g11:func:function}{functies} van $p$, en zoek voor welke
$p$ de \omterm{def:g10:proba:distribution}{kans} om alle drie te scoren gelijk is aan $\frac{27}{64}$.
\end{exercise}

\begin{exercise}[$\star\star$]\label{exo:g11:binom:9}
Hoe vaak moet je met een eerlijke munt gooien opdat de \omterm{def:g10:proba:distribution}{kans} op minstens één
keer kop groter wordt dan $0.99$? (\omterm{def:g10:proba:operations}{Complement}, en probeer daarna
opeenvolgende waarden van $n$.)
\end{exercise}

\begin{exercise}[$\star\star$]\label{exo:g11:binom:10}
Bewijs met de regel van Pascal (\cref{prop:g11:binom:pascal}) en
$\binom n0 = \binom nn = 1$ dat de getallen van elke rij van de \omterm{prop:g11:binom:pascal}{driehoek van
Pascal} optellen tot $2^n$: duid beide leden als een telling van alle paden
van het \omterm{met:g10:proba:tree}{boomdiagram}.
\end{exercise}

\begin{exercise}[$\star\star\star$]\label{exo:g11:binom:11}
Een politicus beweert $60\%$ steun te hebben. In een willekeurige \omterm{def:g10:proba:sample}{steekproef}
van $10$ mensen steunen er slechts $3$ hem.
\begin{enumerate}
  \item Welke verdeling volgt het aantal $X$ steunbetuigingen in de
        \omterm{def:g10:proba:sample}{steekproef} onder die bewering? Bereken $\P(X \leq 3)$.
  \item Is de waarneming verenigbaar met de bewering, volgens de
        beslissingsregel van \cref{met:g11:binom:decision} met een drempel van
        $5\%$?
\end{enumerate}
\end{exercise}

\section{Opgave: de bordplank van Galton}

\begin{problem}[{Weekendopgave --- balletjes, pinnen en de driehoek van
Pascal: hoe de klokvorm geboren wordt, waarom een reeks finales de sterkste
ploeg bevoordeelt, en wanneer je alarm moet slaan}]
\label{pb:g11:binom:1}
Laat duizend balletjes door een rooster van pinnen vallen, waarbij elke
botsing een eerlijke links-of-rechtsmuntworp is, en de vakjes eronder vullen
zich tot een gladde, symmetrische klok --- elke keer opnieuw. Het toestel
heet de bordplank van Galton, en haar wiskunde is precies de binomiale
verdeling van dit hoofdstuk (\cref{thm:g11:binom:binomial}). Deze opgave
bouwt de driehoek, laat de plank draaien, leidt een reeks van zeven finales
in goede banen, en eindigt waar de binomiale verdeling haar loon verdient:
bij de beslissing wanneer een waarneming ons aan een bewering moet doen
twijfelen.

\medskip
\noindent\textbf{Deel I --- De driehoek.}
\begin{enumerate}
  \item Bouw de \omterm{prop:g11:binom:pascal}{driehoek van Pascal} tot rij $6$
        (\cref{prop:g11:binom:pascal}). Formuleer en verklaar de symmetrie
        $\binom nk = \binom{n}{n-k}$ in één zin ($k$ voorwerpen kiezen is
        hetzelfde als\,\dots).
  \item Ga op de rijen $4$ en $5$ na dat elke \omterm{def:g11:seq:sequence}{rij} optelt tot $2^n$, en bewijs
        het: wat tellen alle $\binom nk$ samen?
  \item Leid de regel van Pascal $\binom{n+1}{k} = \binom nk +
        \binom{n}{k-1}$ opnieuw af met het commissieargument: kies één
        bijzondere persoon en splits de commissies op naar diens lot.
  \item Bereken $\binom73$ twee keer: uit de driehoek, en met de
        faculteitsformule.
  \item Ga de trapidentiteit
        $\binom22 + \binom32 + \binom42 + \binom52 = \binom63$ na, en verklaar
        ze door de regel van Pascal vanaf $\binom63$ naar beneden te laten
        cascaderen.
\end{enumerate}

\medskip
\noindent\textbf{Deel II --- De plank.}
Een balletje valt door $n$ \omterm{def:g11:seq:sequence}{rijen} pinnen; bij elke pin stuitert het
onafhankelijk met \omterm{def:g10:proba:distribution}{kans} $\frac12$ naar links of naar rechts. Nummer de vakjes
$0$ tot $n$ volgens het aantal stuiten naar rechts.
\begin{enumerate}[resume]
  \item Leg met de checklist van \cref{met:g11:binom:recognize} uit waarom het
        vakjesnummer de binomiale verdeling
        $\mathcal B\!\left(n, \frac12\right)$ volgt.
  \item Voor een kleine plank ($n = 4$): geef de vijf kansen per vakje. Welk
        vakje raakt het drukst bezet?
  \item Nu $n = 10$ en $1\,024$ balletjes: hoeveel balletjes \emph{verwacht}
        je in het middelste vakje, in vakje $7$ en in elk randvakje?
        Beschrijf de vorm van de stapel.
  \item Bereken voor $X \sim \mathcal B\!\left(10, \frac12\right)$ de waarden
        $\E(X)$, $V(X)$ en $\sigma$ (\cref{prop:g11:binom:expectation});
        bereken daarna het verwachte aandeel balletjes binnen $2\sigma$ van
        het \omterm{prop:g10:coordgeom:midpoint}{midden} (de vakjes $2$ tot $8$) en vergelijk met de garantie van
        Chebyshev uit \cref{pb:g11:stat:1}.
  \item Een scheefgezette plank stuitert met \omterm{def:g10:proba:distribution}{kans} $0.6$ naar rechts: geef
        $\E$, $V$ en $\sigma$ voor $n = 10$, en beschrijf wat er met de stapel
        gebeurt.
  \item In één of twee zinnen: wat in het ontwerp van de plank brengt de
        klokvorm voort --- en waarom stapelen zoveel grootheden uit de
        werkelijkheid (lichaamslengten, meetfouten) zich op dezelfde manier
        op? (De diepe stelling achter beide is de centrale limietstelling, het
        \omterm{pb:g11:scal:1}{hoogtepunt} van de kansrekening in de universitaire volumes.)
\end{enumerate}

\medskip
\noindent\textbf{Deel III --- Best of zeven.}
Twee ploegen spelen een reeks: wie het eerst $4$ wedstrijden wint, pakt de
titel; de wedstrijden zijn onafhankelijk.
\begin{enumerate}[resume]
  \item Gelijkwaardige ploegen ($p = \frac12$): bereken de \omterm{def:g10:proba:distribution}{kans} dat de reeks
        met een schoonveeg eindigt (precies $4$ wedstrijden).
  \item Bereken de \omterm{def:g10:proba:distribution}{kans} dat de reeks de volle $7$ wedstrijden duurt (wat moet
        de stand na $6$ wedstrijden zijn?).
  \item Vervolledig de \omterm{def:g11:prob:rv}{kansverdeling} van de lengte van de reeks ($4$, $5$,
        $6$ of $7$ wedstrijden) voor gelijkwaardige ploegen, en bereken de
        verwachte lengte. Welke lengten zijn het waarschijnlijkst?
  \item Nu wint één ploeg elke wedstrijd met $p = 0.6$. Bereken haar \omterm{def:g10:proba:distribution}{kans} om
        de reeks te winnen (winst in $4$, $5$, $6$ of $7$ wedstrijden: telkens
        wint de ploeg de laatste wedstrijd en $3$ van de vorige). Wat deed de
        reeks met haar voorsprong per wedstrijd?
  \item Vergelijk met één enkele finale ($60\,\%$) en met een reeks van drie
        (bereken die). Formuleer het algemene effect van de lengte van een
        reeks op kunde tegenover geluk --- en waarom competities lange finales
        verkiezen.
\end{enumerate}

\medskip
\noindent\textbf{Deel IV --- Wanneer alarm slaan?}
\begin{enumerate}[resume]
  \item Er wordt $100$ keer met een munt gegooid, met $62$ keer kop. Geef
        voor een eerlijke munt $\E$, $\sigma$ en de \omterm{pb:g11:stat:1}{z-score}
        (\cref{pb:g11:stat:1}) van de waarneming. Wat is het oordeel volgens
        de afspraak van $2\sigma$?
  \item Een leverancier beweert $2\,\%$ defecte stukken te leveren. In een
        partij van $50$ vind je er $3$ defect. Bereken $\P(X \geq 3)$ onder
        die bewering ($X \sim \mathcal B(50, 0.02)$; ga via $\P(X = 0)$,
        $\P(X = 1)$ en $\P(X = 2)$). Alarmerend bij de drempel van $5\,\%$
        (\cref{met:g11:binom:decision}, \cref{exo:g11:binom:11})?
  \item Volhouden bij de loterij: elk lot wint (iets) met \omterm{def:g10:proba:distribution}{kans}
        $\frac{1}{1000}$. Bereken de \omterm{def:g10:proba:distribution}{kans} op minstens één winnend lot bij
        $1\,000$ loten. Het antwoord ($\approx 63\,\%$, geen $100\,\%$!)
        verbergt een beroemde constante: bereken $0.999^{1000}$ en houd het
        getal $0.368$ in gedachten voor jaar 12.
  \item Slotstuk --- het portret van de binomiale verdeling: de checklist om
        haar te herkennen (vaste $n$, onafhankelijkheid, constante $p$); de
        \omterm{prop:g11:binom:pascal}{driehoek van Pascal} als haar tabel; de klok als haar vorm; $np$ en
        $np(1-p)$ als haar kompas; en haar twee erfgenamen die in jaar 12
        wachten --- de gladde klokkromme en de wet van de grote aantallen.
        Telkens één zin.
\end{enumerate}
\end{problem}
